% Danh mục từ viết tắt (đặt sau Danh mục bảng biểu theo hướng dẫn Khoa HTTT).
\chapter*{Danh mục từ viết tắt}
\addcontentsline{toc}{chapter}{Danh mục từ viết tắt}

\begin{longtable}{@{}p{3.2cm}p{11.5cm}@{}}
\toprule
\textbf{Từ viết tắt} & \textbf{Ý nghĩa} \\
\midrule
\endhead
\bottomrule
\endfoot
API     & Giao diện lập trình ứng dụng (\emph{Application Programming Interface}) \\
AP      & Độ chính xác trung bình (\emph{Average Precision}) \\
AUC     & Diện tích dưới đường cong (\emph{Area Under the Curve}) \\
CSKH    & Chăm sóc khách hàng \\
CSV     & Tệp dữ liệu phân tách bằng dấu phẩy (\emph{Comma-Separated Values}) \\
EDA     & Phân tích khám phá dữ liệu (\emph{Exploratory Data Analysis}) \\
F1      & Trung bình điều hòa của precision và recall \\
HHI     & Chỉ số Herfindahl--Hirschman \\
JWT     & Chuỗi xác thực dạng JSON (\emph{JSON Web Token}) \\
ML      & Học máy (\emph{Machine Learning}) \\
PR      & Precision--Recall \\
PR-AUC  & Diện tích dưới đường cong precision--recall \\
RF      & Rừng ngẫu nhiên (\emph{Random Forest}) \\
ROC     & Đường cong đặc tính hoạt động của bộ thu (\emph{Receiver Operating Characteristic}) \\
ROC-AUC & Diện tích dưới đường cong ROC \\
SHAP    & Giá trị giải thích dựa trên lý thuyết trò chơi (\emph{SHapley Additive exPlanations}) \\
SLA     & Cam kết thời hạn xử lý (\emph{Service Level Agreement}) \\
SMOTE   & Kỹ thuật sinh mẫu tổng hợp cho lớp thiểu số (\emph{Synthetic Minority Over-sampling Technique}) \\
SVM     & Máy vectơ hỗ trợ (\emph{Support Vector Machine}) \\
\end{longtable}
